\chapter{Web architecture, REST, and security boundaries}
\chapterintro{A useful HTTP server lives in an ecosystem of TLS terminators, proxies, caches, clients, and standards. This chapter completes the theory gate and defines exactly what we will build.}

\begin{objectives}
\begin{itemize}
  \item Explain forward proxies, reverse proxies, gateways, and caches.
  \item Describe TLS goals and the trust decisions behind HTTPS.
  \item Apply REST constraints to a resource-oriented API without treating REST as a synonym for JSON.
  \item Compare HTTP/1.1, HTTP/2, and HTTP/3 at a conceptual level.
  \item State the supported, rejected, and explanation-only features of the project.
\end{itemize}
\end{objectives}

\section{Intermediaries are normal}

HTTP was designed to pass through intermediaries. A \term{forward proxy} is
selected on behalf of a client. A \term{reverse proxy} accepts traffic on behalf
of origin applications. A \term{gateway} translates or fronts another system.
A \term{cache} reuses stored responses when protocol rules allow.

A reverse proxy may terminate TLS, enforce request limits, authenticate users,
balance traffic, compress responses, and forward a simpler request to an
application server. This is a common production boundary. It does not excuse
the application from validating inputs, and it creates a trust question: fields
such as a forwarded client address are meaningful only when set or sanitized by
a trusted intermediary.

\section{Caching changes the path}

A cache can answer a request without contacting the origin. Cache keys, freshness
lifetimes, validators, request directives, response directives, and method/status
rules determine whether reuse is valid. An \code{ETag} validator can support a
conditional GET; if the stored representation is still current, the origin can
return \status{304 Not Modified} without the representation content.

Caching is not implemented in \code{tinyhttp}, but it is important to API design.
Unsafe methods should invalidate relevant stored responses. Personalized data
must not accidentally become shared. A backend developer must understand that a
response can affect later requests even after its connection closes.

\section{What TLS provides}

TLS wraps a transport connection to provide confidentiality, integrity, and
authenticated server identity when certificates are validated correctly. During
the handshake, peers negotiate cryptographic parameters, the server presents a
certificate chain, and the client checks trust, hostname, validity, and policy.
The handshake establishes keys used to protect application data.

TLS does not make a vulnerable application safe. It protects data in transit
between TLS peers. An authorized client can still send malicious input; an
application can still expose secrets in a response; decrypted data still exists
at endpoints. Certificate verification is essential---encryption to an
unverified attacker is not the intended security property of HTTPS.

Our learning server does not implement cryptography and does not terminate TLS.
Production deployments normally place an established TLS implementation and
production server or proxy in front of application code.

\section{REST as an architectural style}

REST is not ``HTTP plus JSON'' and not a route naming convention. Its constraints
include client--server separation, stateless requests, cacheability, a uniform
interface, a layered system, and optional code-on-demand. The uniform interface
includes identification of resources, manipulation through representations,
self-descriptive messages, and hypermedia as the engine of application state.

Many practical APIs use several REST ideas without implementing mature
hypermedia. This book will be precise: the Task API is resource-oriented and
uses HTTP semantics carefully, but it does not claim every possible REST
maturity property.

\subsection{Model nouns before endpoints}

The core collection resource is \code{/tasks}; an individual task is
\code{/tasks/\{id\}}. Methods carry operations, so paths do not become
\code{/createTask} or \code{/deleteTask}. Representations contain task fields.
Filters and pagination refine the collection request rather than inventing a
new verb-shaped endpoint.

Statelessness means every request carries what the server needs to interpret it.
Database state persists tasks, but the meaning of PATCH does not depend on a
secret sequence of prior requests on the same TCP connection.

\section{Input is hostile until proven otherwise}

Security is an end-to-end property. The server will enforce:

\begin{itemize}
  \item strict message syntax and unambiguous framing;
  \item finite line, field, body, time, worker, and queue limits;
  \item header-value validation to prevent response splitting;
  \item media-type and UTF-8 checks before JSON parsing;
  \item allow-lists, lengths, types, and formats for task fields;
  \item parameterized SQL rather than string interpolation;
  \item generic 500 responses while logging internal exceptions;
  \item loopback binding by default.
\end{itemize}

Authentication is intentionally absent. A toy password or static token can
teach unsafe habits if mistaken for a deployable identity system. The security
chapter explains where authentication and authorization belong and why they
need established password hashing, token/session policy, secret management,
revocation, TLS, and audit design.

\section{Why newer HTTP versions exist}

HTTP/1.1 can reuse connections but concurrent requests often require several
connections, and ordered responses can delay later work. HTTP/2 retains HTTP
semantics while using binary frames and multiplexed streams on one connection.
It compresses fields using HPACK. Packet loss still affects the underlying TCP
connection, even if application streams are independent at the HTTP layer.

HTTP/3 maps HTTP semantics to QUIC, which runs over UDP and provides encrypted,
multiplexed streams with loss recovery. Loss affecting one QUIC stream need not
block delivery on unrelated streams in the same way TCP ordering can. QPACK
handles field compression under these constraints.

We implement HTTP/1.1 because a human can inspect its wire messages and because
it makes framing visible. We do not pretend the same parser can be adapted to
HTTP/2 by adding a condition; the binary framing and connection machinery are
different implementations of shared semantics.

\section{The implementation contract}

\begin{center}
\begin{tabularx}{\textwidth}{@{}>{\bfseries}p{30mm} X X@{}}
\toprule
Category & Included & Boundary \\
\midrule
Protocol & HTTP/1.1 origin-form, Host, Content-Length, core methods, HEAD/OPTIONS, persistent sequential requests & Other versions rejected \\
Framework & request/response types, multivalue fields, router, middleware, problem responses & No templates, static files, cookies, sessions, or dependency injection \\
Transport & IPv4 loopback TCP, bounded blocking workers, timeouts, clean shutdown & No TLS, IPv6 listener, event loop, or production tuning \\
Content & UTF-8 JSON with 1 MiB request limit & Chunking, streaming, compression, ranges, and negotiation explained but not implemented \\
Backend & SQLite Task CRUD, validation, filtering, pagination, transactions & No accounts, authentication, or distributed database \\
\bottomrule
\end{tabularx}
\end{center}

An explicit boundary makes the project honest. Unsupported features receive a
clear error rather than accidental behavior. The final production bridge will
show how WSGI, ASGI, mature servers, reverse proxies, and frameworks address
broader needs.

\begin{checkpoint}
Before Part II, you should be able to explain, without code, how a URL becomes
an HTTP request; why TCP has no request boundary; how HTTP/1.1 finds the end of
a length-framed body; why methods and status codes have semantics; and which
features our strict subset rejects. If one explanation is vague, revisit the
relevant chapter before implementing its parser.
\end{checkpoint}

\summarytitle

HTTP systems include trusted and untrusted intermediaries, caches, TLS, and
multiple protocol versions. REST supplies architectural constraints, not a JSON
branding label. The server we build has a deliberately strict contract: enough
HTTP/1.1 for a meaningful backend, explicit limits, and clear production
boundaries.

\reviewtitle
\begin{enumerate}
  \item When may a reverse proxy field describing client identity be trusted?
  \item Why does TLS not replace application input validation?
  \item Which REST constraint is violated by hidden per-connection application state?
  \item Why is HTTP/2 not simply HTTP/1.1 with more headers?
\end{enumerate}

\begin{exercise}[Theory gate]
On one page, diagram a POST from a client to the future Task API through DNS,
TCP, an optional TLS proxy, HTTP parsing, routing, validation, SQLite, and the
response. Label every byte/text conversion and every trust boundary.
\end{exercise}

\begin{exercise}[Feature classification]
Classify each as implemented, explicitly rejected, or explanation-only:
\code{Content-Length}, chunked request bodies, GET, HTTP/2, TLS, repeated ordinary
fields, duplicate Content-Length, JSON, authentication, and keep-alive. Explain
why each boundary is visible rather than silent.
\end{exercise}

